\section{Introduction}
\label{sec:intro}

Algorithmic redistricting proposes a simple solution to partisan gerrymandering:
replace human judgment with a mathematical algorithm that optimises a
politically neutral criterion.
The leading candidate --- minimum edge-cut redistricting --- minimises the
total length of district boundaries subject to equal-population constraints,
producing compact districts without any reference to which party's voters
live where \citep{karypis1998,b8paper,b9paper}.

The algorithm is genuinely neutral.
Yet its outcomes are not.
In Wisconsin, a 50.3\,\% Democratic state, minimum edge-cut redistricting
consistently produces 3 Democratic seats out of 8, a proportionality gap of
$-12.8$ percentage points.
In Georgia, a near-even 50.1\,\% Democratic state, the algorithm produces
3 of 14 Democratic seats ($-14.4$\,pp).
These results are not idiosyncratic --- across all 8 competitive states
(46--55\,\% Democratic vote share) in our 50-state sweep, every single one
shows Republican over-representation, with a mean gap of $-6.2$\,pp
\citep{b9paper}.

\citet{rodden2019} provides the explanation: Democratic voters are geographically
concentrated in dense urban cores, while Republican voters are spread across
suburban and rural areas.
Any compact-district algorithm --- because it minimises boundary length ---
tends to draw boundaries that respect this geographic sorting, producing
implicit Democratic packing without any partisan input.
The Rodden effect is a fact about geography, not about algorithms.

This raises a precise mathematical question: \emph{what partisan information
is required to guarantee proportional outcomes, and how far do geography-only
algorithms fall short?}
We call the difference the \emph{proportionality gap} and characterise it
across all 50 states.

\paragraph{This paper's contribution.}
We introduce a dual-constraint bisection framework that extends minimum-edge-cut
redistricting to METIS's $\mathtt{ncon}=2$ mode with vertex weights
$[\text{population},\, D_{\text{votes}}]$.
The Democratic-vote constraint encodes how many Democratic voters belong to
each half of a bisection.
By choosing these targets precisely, we can guarantee that each half produces
the correct number of party-winning districts.

The central result is a closed-form formula (Theorem~\ref{thm:proportionality})
for the $\mathtt{tpwgts}$ parameters that achieve exact partisan proportionality.
The formula is derived from first principles: the R-bloc receives the minimum
Democratic voter share needed for Republicans to win every district in that bloc,
and no more.

The companion result is a feasibility characterisation (Theorem~\ref{thm:lorenz-feasibility}):
a \emph{partisan Lorenz curve}, sorting tracts from most to least Democratic,
determines when proportional packing is achievable under the geographic
constraint.
States with uniform partisan geography (near-diagonal Lorenz curves) can
achieve proportionality easily.
States with extreme urban concentration (Lorenz curves far above the diagonal)
cannot achieve proportionality through any contiguous redistricting.

\paragraph{What this paper is not.}
We do not advocate for partisan redistricting.
The Districting Integrity Act's \S104(e) prohibition on partisan inputs for
federal congressional redistricting reflects a considered policy judgment that
we do not contest.
The calculation we perform is analogous to computing, for a chemist, the energy
required to break a bond --- it measures a physical quantity without advocating
for the reaction.
Specifically, we compute the \emph{minimum partisan signal} required to close
the proportionality gap, and characterise when that signal is zero (when
geography-only algorithms happen to produce proportional outcomes) versus
large (when geographic concentration makes proportionality geometrically expensive).

\paragraph{Paper structure.}
Section~\ref{sec:related} reviews related work on efficiency gaps, partisan
Lorenz curves, and dual-constraint graph partitioning.
Section~\ref{sec:framework} introduces the $\mathtt{ncon}=2$ framework and
derives the $\mathtt{tpwgts}$ formula.
Section~\ref{sec:theory} proves the proportionality guarantee and the
Lorenz feasibility condition.
Section~\ref{sec:eval} presents the 50-state sweep.
Section~\ref{sec:compromise} characterises states where proportionality is
cheaply achievable.
Section~\ref{sec:legal} discusses legal implications.
Section~\ref{sec:conclusion} concludes.
